
%\documentclass[landscape]{slides}
\documentclass{beamer}
\usepackage{beamerfoils}
\usetheme{Pittsburgh}
%\usepackage[accumulated]{beamerseminar}
                                % remove ``accumulated'' option
                                % for original behaviour
%\usepackage{beamerthemeclassic}
\usepackage{ulem}
%\usepackage{beamerthemesplit}

%\usepackage{beamerthemeshadow}

%\usepackage{beamerthemeplain}
%\usepackage{times}
\usepackage{colortbl}

%\usepackage{lslide}
\def\emptyline{\vspace{10pt}}
\usepackage{amssymb}
%\usepackage{french}
%\usepackage[latin1]{inputenc}
\usepackage[utf8]{inputenc}
%\usepackage[T1]{fontenc}
\usepackage[]{color}
%\usepackage[usenames,dvipsnames]{color}
%\usepackage[pdftex]{color}
%\usepackage[]{colorlslide,epsfig}
%\usepackage{amsfonts,stmaryrd}
\usepackage{epsfig}

\usepackage{graphicx,array}
\usepackage{tikz}
\usetikzlibrary{shadows.blur}
\usepackage{amsmath}%,amsthm}
\usepackage{hyperref}
\usepackage{pifont}

%\renewcommand{\frame}{\end{slide}   \begin{slide}}
%\newcommand{\frametitle}{\centerline}
 

%\usepackage{beamerthemeshadow}
\usepackage{graphicx}
%\usepackage{beamerthemesplit}



%%%%%%%%%%%%%%%%%% ENVIRONEMENT %%%%%%%%%%%%%%%%%%%%%%%%%%%%%%%%%%%%%%%%%%%%%%%%%
\newtheorem{deft}{Définition}
%\newtheorem{theorem}{Théorème}
\newtheorem{proposition}{Proposition}
\newtheorem{exemple}{Exemple}
\newtheorem{postulat}{Postulat}


%%%%%%%%%%%%%%%%%% COULEURS %%%%%%%%%%%%%%%%%%%%%%%%%%%%%%%%%%%%%%%%%%%%%%%%%
%\definecolor{coltitle}{rgb}{0.1, 0.1, 0.543}
%\definecolor{colitemi}{rgb}{0, 0.75, 0.99}
%\definecolor{colexemp}{rgb}{0.6, 0.4, 0}
%\definecolor{colmult}{rgb}{0, 1, 0}
%\definecolor{gris}{rgb}{0.8, 0.8, 0.8}
\definecolor{encre}{rgb}{0.2, 0, 0.4}
\definecolor{lightblue}{rgb}{0.3,0.35,5}
% Couleurs pour le diagramme de classification d'Adleman
\definecolor{adlemanblue}{RGB}{30,30,120}
\definecolor{adlemangray}{RGB}{192,192,192}
%\definecolor{grisbleu}{rgb}{0.8, 0.8, 1}
%\definecolor{cadre}{rgb}{0.7, 0.7, 0.7}
%\definecolor{cadresubsec}{rgb}{0, 0.3, 0.3}
%\definecolor{cadreombre}{rgb}{0.8, 0.8, 0.8}
%\definecolor{colsection}{rgb}{0.4, 0.2, 0.2}

%\definecolor{bleur}{rgb}{0.3, 0, 0.6}
%\definecolor{bleuv}{rgb}{0, 0.3, 0.6}
%\definecolor{rougeb}{rgb}{0.6, 0, 0.3}
%\definecolor{rougev}{rgb}{0.6, 0.3, 0}
%\definecolor{rougealert}{rgb}{1., 0.2, 0}
%\definecolor{vertr}{rgb}{0.3, 0.6, 0}
%\definecolor{vertb}{rgb}{0, 0.6, 0.3}
%\definecolor{blanc}{rgb}{1, 1, 1}

\newcommand{\red}[1]{\only<#1>{\color{red}}}
\newcommand{\green}[1]{\only<#1>{\color{green}}}
\newcommand{\lra}{\leftrightarrow}
\newcommand{\Fd}{\mathbb{F}}
\newcommand{\enlight}[1]{\textit{\color{blue}#1}}
\newcommand{\blue}[1]{\textit{\color{blue}#1}}

%%   for the headline etc.
\renewcommand\normalcolor{\color{coltitle}}


%%   for the background
%\pagecolor{grisbleu}
%%   for the text
\color{encre}

%%%%%%%%%%%%%%% voir le style ....
 \beamertemplatenavigationsymbolsempty
\beamertemplatefootempty
%\beamertemplatefootframenumber
\setbeamertemplate{footline}[frame number]


\title{Logiciels Malveillants - Généralités}
\date{Février 2026}

\AtBeginSection[]
{
\begin{frame}
\frametitle{Plan}
\tableofcontents[currentsection,sectionstyle=show/shaded,subsectionstyle=show/show/hide]
\end{frame}
}

\AtBeginSubsection[]
{
\begin{frame}
\frametitle{Plan}
\tableofcontents[currentsubsection,sectionstyle=show/shaded,subsectionstyle=show/shaded/hide]
\end{frame}
}
%%%%%%%%%%%%%%%%%%%%%%%%%%%%%%%%%%%%%%%%%%%%%%%%%%%%%%%%%%%%%%%%%%%%%%%%%%%%%%%%%%
\titlegraphic{
\centering

\includegraphics[width=4cm]{logo-upvd.png}\\[0.3cm]

\includegraphics[width=6cm]{malware1.jpg}
}
\begin{document}

\maketitle


\section{Introduction - Taxonomie}

\subsection{Definition et classification}



\begin{frame}
\frametitle{Définition d'une infection informatique}

%\begin{deft}[Infection informatique]
 \begin{block}{Définition (Infection informatique) :} 
  Un \enlight{programme simple} ou \enlight{auto reproducteur} :
\begin{itemize}
\item  à caractère offensif, s'installant dans un système d'information, à l'insu du ou des utilisateurs,
\item en vue de porter atteinte à la confidentialité, l'intégrité ou la disponibilité de ce système,
  \item ou susceptible d'incriminer à tort son possesseur ou l'utilisateur dans la réalisation d'un crime ou d'un délit.
\end{itemize}
\end{block}
%\end{deft}

\end{frame}



\begin{frame}
  
\begin{postulat}
Les infections informatiques concernent tous les environnements capables d'exécution !
\end{postulat}

\begin{itemize}
\item Tous les systèmes d'exploitation.
\item Les environnements mobiles (téléphones, consoles de jeux, GPS, ordinateurs embarqués \ldots).
\item De nombreux formats de documents.
\end{itemize}
\end{frame}

\begin{frame}
\frametitle{Classification d'Adleman}

Adleman a classifié les différents agents infectieux informatiques en 4 catégories.

\begin{center}
\begin{tikzpicture}[
    bluebox/.style={rectangle, rounded corners=10pt, fill=adlemanblue, text=white,
                    minimum height=1cm, minimum width=2.4cm, align=center, font=\footnotesize\bfseries,
                    blur shadow={shadow blur steps=5, shadow xshift=0.5mm, shadow yshift=-0.5mm}},
    graybox/.style={rectangle, rounded corners=10pt, fill=adlemangray, text=black,
                    minimum height=1cm, minimum width=3cm, align=center, font=\footnotesize\bfseries,
                    blur shadow={shadow blur steps=5, shadow xshift=0.5mm, shadow yshift=-0.5mm}},
    connector/.style={line width=1.8pt, rounded corners=8pt, draw=black!70}
]
    % Niveau 1 - Racine
    \node[bluebox, minimum width=3.5cm] (malware) at (0,0) {Infections\\informatiques};

    % Niveau 2
    \node[graybox] (simple) at (-3,-2) {Programmes\\simples};
    \node[graybox, minimum width=3.5cm] (selfrep) at (3,-2) {Programmes\\autoreproducteurs};

    % Niveau 3
    \node[bluebox] (logic) at (-4.3,-4) {Bombes\\logiques};
    \node[bluebox] (trojan) at (-1.6,-4) {Chevaux\\de Troie};
    \node[bluebox] (virus) at (1.6,-4) {Virus};
    \node[bluebox] (worm) at (4.3,-4) {Vers};

    % Connexions avec coins arrondis
    \draw[connector] (malware.south) -- ++(0,-0.5) -| (simple.north);
    \draw[connector] (malware.south) -- ++(0,-0.5) -| (selfrep.north);
    \draw[connector] (simple.south) -- ++(0,-0.5) -| (logic.north);
    \draw[connector] (simple.south) -- ++(0,-0.5) -| (trojan.north);
    \draw[connector] (selfrep.south) -- ++(0,-0.5) -| (virus.north);
    \draw[connector] (selfrep.south) -- ++(0,-0.5) -| (worm.north);
\end{tikzpicture}
\end{center}

La différence majeure est
\begin{itemize}
\item \enlight{Programme simple} : programme classique ayant un objectif malicieux.
\item \enlight{Auto-reproducteur} : programme qui recopie son propre code, en général, dans d'autres exécutables.
\end{itemize}
\end{frame}

\subsection{Bombe logique}

\begin{frame}
\frametitle{Bombe logique}

\begin{deft}[Bombe Logique]
Une bombe logique est un programme infectant simple s'installant dans le système et qui attend un évènement (date, action, données particulières...) appelé en général \blue{gâchette} pour exécuter sa fonction offensive. 
\end{deft}


Ces programmes constituent souvent la charge finale d'un virus (exp. virus CIH qui se déclenche chaque 26 avril)
\end{frame}


\begin{frame}
\frametitle{Exemple de bombe logique}
\begin{itemize}
\item Un administrateur système avait implanté un programme vérifiant la présence de son nom dans les registres de feuilles de paye de son entreprise.

\item En cas d'absence de ce nom (signifiant un licenciement) le programme chiffrait tous les disques dur de l'entreprise.
  \begin{center}
    \includegraphics[width=5cm]{encrypted-hard-drive.jpeg}
    \end{center}
\item L'entreprise ne possédant pas la clef de chiffrement ne pouvait plus accéder à ses données. 
\end{itemize}
\end{frame}

\begin{frame}
\frametitle{Limite de la lutte anti-virale contre les bombes logiques}
\begin{itemize}
\item Les anti-virus ont du mal à lutter contre les bombes logiques, ce sont en  apparence de simples programmes.
\item Les techniques évoluées de lutte anti-virale (analyse euristique, émulation de code) sont condamnées à être constamment mises en défaut face à des bombes logiques inconnues.
\end{itemize}
\end{frame}


\subsection{Chevaux de Troie}

\begin{frame}
\frametitle{Chevaux de Troie}

\begin{deft}[Cheval de Troie]
Un cheval de Troie est un programme simple, composé de deux parties : 
\begin{itemize}
\item le module serveur :
  \begin{itemize}
  \item Il est installé dans l'ordinateur de la victime.
    \item  donne discrètement à l'attaquant l'accès à tout ou partie des ressources via le réseau.
    \end{itemize}
\item  le module client :
  \begin{itemize}
  \item  il se trouve sur la machine de l'attaquant
    \item il permet d'accéder aux services délivrés inconsciemment par la victime.
    \end{itemize}
\end{itemize}




\end{deft}

\end{frame}



\begin{frame}
\frametitle{Chevaux de Troie}
{\color{blue} Le module serveur :}
\begin{itemize}
\item généralement dissimulé dans un autre programme \textit{\color{blue} anodin et attractif.}
  \item L'exécution de ce dernier  installe à l'insu de la victime la partie serveur du cheval de Troie.
\end{itemize}
\pause 
{\color{blue} Le module client :}
\begin{itemize}
\item recherche sur le réseau  l'adresse IP et le port (TCP ou UDP) les machines infectées par le module serveur.
  \item prend le contrôle des machines infectées trouvées. 
%\item Cette prise de contrôle lui permet de mener des actions offensives.
\end{itemize}
\end{frame}


\begin{frame}
\frametitle{Client-Serveur dans les chevaux de Troie}

\includegraphics[width=10cm]{client-serveur}

\end{frame}

\begin{frame}
\frametitle{Détection des chevaux de Troie}

Ces  programmes sont détectés de manière inégale.
\begin{itemize}
\item Un cheval de Troie non diffusé sur internet, bien programmé a toutes les chances de contourner un anti-virus.
  \bigskip 
\item L'usage d'un pare-feu est hautement conseillé, même si plusieurs techniques sont connues ou à l'étude qui permettent de les contourner.
\end{itemize}
\end{frame}

\begin{frame}
\frametitle{Exemple de chevaux de Troie}
\begin{itemize}
\item Les leurres :  programmes imitant le fonctionnement normal d'un programme légitime du système (fausse bannière de connexion Unix par exemple)
  \bigskip
\item Les espions de clavier, \enlight{keyloggers}, ne sont que des cas particuliers de chevaux de Troie. 
\end{itemize}
\bigskip
L'action offensive consiste le plus souvent à récupérer une ou plusieurs informations présent sur la machine de la victime.
\end{frame}





\subsection{Virus et Ver}
\begin{frame}
\frametitle{Ver et Virus : définitions }
\begin{deft}[Virus] Programme se propageant par recopie de son propre code.
\end{deft}

\begin{deft}[Ver] Peut être vu comme un virus orienté réseau. La différence éventuelle avec un virus réside dans le fait qu'un ver n'est pas forcément attaché à un fichier (ver simple de type slammer).
\end{deft}
\end{frame}

\begin{frame}
\frametitle{Les vers informatiques}

Trois classes principales de vers
\begin{itemize}
\item {\color{lightblue} Les vers simples.} Utilisations d'une faille de sécurité (Slammer, Sasser,\ldots).
\item {\color{lightblue} Les macros vers.} Utilisations d'ingénieurie sociale et d'une pièce jointe bureautique infectée.
\item {\color{lightblue} Les vers de courrier électronique.} Utilisation d'ingénieurie sociale et d'une pièce jointe exécutable infectée (Bagle, Netsky,\ldots).
\end{itemize}

\end{frame}


\begin{frame}
Par la suite nous nous limiterons essentiellement au programme auto-reproducteur:
\begin{itemize}
\item nous étudierons essentiellement les virus,
\item et dans une moindre mesure les vers.
\end{itemize}
\end{frame}

\section{Généralité sur les virus}

\subsection{Mode de propagation et cycle de vie}

\begin{frame}
\frametitle{Mode de propagation}
%\begin{enumerate}
% \item
  Le programme infectant proprement dit est porté par un programme hôte ou \enlight{dropper}.
  % (dit programme infecté, dans le cas de la première infection réalisé par  l'attaquant, le terme de \enlight{dropper} (i.e. largueur) est utilisé).
  \bigskip
  
  {\color{blue} Lorsque le \enlight{dropper} est exécuté : }
\begin{enumerate}
\item Le programme infectant (ou code viral) prend la main :
  \begin{itemize}
  \item il agit selon son mode propre
  \item le programme hôte est  temporairement mis en sommeil.
  \end{itemize}
  \medskip
  
\item il redonne la main au programme hôte qui s'exécute alors normalement sans trahir la présence du programme infectant.
\end{enumerate}
%\end{enumerate}

\end{frame}


\begin{frame}
\frametitle{Diagramme fonctionnel d'un virus }

\begin{enumerate}
\item Une routine de \enlight{recherche} des programmes cibles
%\begin{itemize}
%\item Vérification que le fichier est exécutable, a le bon format.
%\item Vérification que le fichier n'est pas déjà infecté (\enlight{contrôle de surinfection}).
%\end{itemize}
%Cette partie détermine son efficacité en termes de portée et de rapidité.
\medskip
\item Une routine de \enlight{copie} : injection du code viral dans la cible.
  \medskip
\item Une routine d'\enlight{anti-détection}. Son but: contrer l'action des antivirus.
  \medskip
\item La \enlight{Charge finale} couplée éventuellement à un mécanisme de déclenchement différé (\enlight{gâchette}).% C'est elle qui a un pouvoir de nuisance.
\end{enumerate}
\end{frame}


\begin{frame}
\frametitle{Cycle de vie d'un virus ou d'un ver}

On peut identifier trois phases dans la vie d'un virus
\begin{itemize}
\item La phase d'infection.
\item La phase d'incubation.
\item La phase de maladie.
\end{itemize}

%\enlight{La diffusion} : une fois le virus incorporé à un programme \textit{inoffensif}, le dropper, le programmeur utilisera en fonction de la ou des cibles un programmes adaptés au virus :
%\begin{itemize}
%\item Pour un nombre de victime réduit : nécessité de renseignements préalable. 
%\item Jeux infectés, ou logiciels basés sur l'humour.
%\end{itemize}
\end{frame}


\begin{frame}
\frametitle{Phase d'infection}

Dans cette phase le virus va se propager dans l'environnement
\begin{itemize}
\item \enlight{Passivement}
\begin{itemize}
\item  le support est copié sur un support (CDROM, site de téléchargement, forum de discussion,\ldots) et transmis.
  \medskip 
\item Les victimes peuvent alors le copier et l'exécuter dans leur propre environnement.
\end{itemize}
\medskip

\item \enlight{Activement.}
\begin{itemize}
\item  L'utilisateur exécute le dropper ou un fichier déjà contaminé lors d'une infection antérieure.
  \end{itemize}
\end{itemize}
% C'est cette phase qui différencie les infections simples des virus : la duplication du code. 
\end{frame}

\begin{frame}
\frametitle{Phase d'infection}
Lors de cette phase le virus  procède à la recherche de programmes cibles : %qui seront à leur tour infectés
\begin{itemize}
\item un virus efficace vérifiera que le fichier est exécutable, a le format correct et que la cible n'est pas déjà infectée.
\item il s'assure aussi de minimiser les risques de détections: par exemple une infection d'exécutables *.COM ne doit pas dépasser 64koctets (taille max des fichier *COM).
\end{itemize}

Le virus copie alors son propre code dans la cible.

\end{frame}



\begin{frame}
\frametitle{Phase d'incubation}
\begin{itemize}
\item C'est la plus grande partie de la vie d'un virus. 
\item Mission principale : assurer la survie d'un virus. 
\item Il s'agit de limiter voir d'empêcher sa détection :
\begin{itemize}
\item Une détection par l'utilisateur : en évitant par exemple des erreurs d'exécution.% qui pourraient alerter l'utilisateur.
  \bigskip 
\item Une détection par l'anti-virus : se dérober à la suveillance anti-virale. 

  %Le virus utilisera plusieurs techniques qui vont lui permettre de se dérober à la surveillance anti-virale.
\end{itemize}
\end{itemize}
\end{frame}



\begin{frame}
\frametitle{Phase de maladie}
Lors de cette phase \enlight{la charge finale est activée}, son déclechement dépendra de sa position: %. Son mode de déclenchement dépendra de nombreux facteurs et sera fonction de l'endroit dans le code où la routine sera placée.
\begin{itemize}
\item \enlight{En tête de code},  la charge finalement sera exécutée avant toute infection. %Ce cas est rare, il a pour conséquence de limiter généralement la phase de survie du virus ou du ver.
\item \enlight{En fin de code}, elle n'aura lieu qu'après le processus d'infection.
\item \enlight{Au milieu du code}, en particulier si elle est conditionnée par la réussite ou non de l'infection. 
\end{itemize}
\end{frame}


\begin{frame}
  \frametitle{Phase de maladie : déclenchement}
  \begin{itemize}
\item Le déclenchement peut être différé (gâchette).

\item  La charge finale est alors une bombe logique montée sur un vecteur viral.
  \end{itemize}
  \medskip
  
{\color{blue} Exemples :}


\begin{itemize}
\item Date du système.
\item Nombre d'infections réalisées.
\item La frappe d'une séquence particulière de touches, un certain nombre de fois (e.g. \texttt{CTRL+ALT+SUPP}).
\item Nombre d'ouverture d'un document \texttt{Word}.
\end{itemize}
%En fait tout dépend de l'imagination et des effets recherchés par le programmeur.
\end{frame}





\begin{frame}
\frametitle{Phase de maladie}
L'effet de la charge finale sont de nature très variable :
\begin{itemize}
\item \enlight{de nature non létale} : affichage d'images ou d'animations, de messages, ... % émissions de son ou de musiques (attaque ludique ou militant)
\item \enlight{de nature létale } : il s'agit là de porter atteinte
\begin{itemize}
\item  à la confidentialité des données (évasion des données).
  \smallskip
  
\item  à l'intégrité du système ou des données ( destructions/modifications de fichiers, ...) % modifications aléatoires de données,\ldots)
  \smallskip
  
  \item à la disponibilité du système ( saturation du processeur,  ...)
    \smallskip
    
  \item à la disponibilité des données (chiffrement)  ou incrimination des utilisateurs.% (introduction de données compromettantes, ...))
\end{itemize}
\end{itemize}
\end{frame}


%\end{document}

\begin{frame}
\frametitle{Phase de maladie : quid de la destruction matérielle ?}
\begin{itemize}
\item Généralement non admise par les experts.
\item Controverse du CIH :
  \begin{itemize}
  \item Le CIH ne détruit pas directement le matériels mais les logiciels associés (firmware, etc.).
  \item L'attaque matérielle est seulement simulée.
  \end{itemize}
\item D'autres exemples plus anciens : attaque les lecteurs de disquettes ou disque dur par des requêtes E/S répétées.
%\item Les virus sont alors ciblé sur certains matériels, et sont donc moins détectable par les antivirus.
\end{itemize}
\end{frame}



\begin{frame}
\frametitle{Comparaison biologique/informatique}
{\scriptsize
\begin{center}
\begin{tabular}{|c|c|}
\hline
Virus biologique & Virus informatique \\
\hline 
\hline
Attaques spécifiques de cellules & Attaques spécifiques \\
de format & de format \\
\hline
Les cellules touchées & Le programme infecté \\
produisent de nouveaux  & génère d'autres \\
virus & programmes viraux \\
\hline
Modification  de l'information & Modification des \\
génétique de la cellule & actions du programmes \\
\hline
Le virus utilise les & Multiplication uniquement \\
structures de la  cellule & via un programme infecté \\
hôte pour se multiplier & \\
\hline
Intéractions virales & Virus binaires ou virus anti-virus \\
Multiplication uniquement & Nécessité d'une exécution \\
dans des cellules vivantes & our la dissémination \\
\hline
Une cellule infectée & Lutte contre la surinfection \\
n'est pas suringectée & \\
par le même virus & \\
\hline
Rétrovirus & Virus luttant spécifiquement contre \\
& un anti-virus - Virus de code source \\
\hline
Mutation virale & Polymorphisme viral \\
\hline
Porteur sain & Virus latents \\
\hline
Antigènes & Signatures \\
\hline
\end{tabular}
\end{center}
}
\end{frame}





\subsection{Données et indices numériques}


\begin{frame}
\frametitle{Quelques chiffres}
Les statistiques concernant les malwares sont difficiles à obtenir et à vérifier (...)

%\begin{itemize}
%\item Le nombre total de virus connu était d'environ 70 000 en janvier 2002.
%\item Chaque mois entre 800 et 12000 nouveaux virus sont découverts.
%\item En janvier la répartition entre les différents types d'infection est donnée ci-dessous.
%\end{itemize}


%$$diagramme$$
\begin{center}
  %\includegraphics[width=4cm]{repartition-infection}
  \includegraphics[width=6cm]{progression-nb-malware.jpg} \vspace{0.2cm}
  \includegraphics[width=4.5cm]{repartition-malware-v1.png}
\end{center}


\end{frame}


%\begin{frame}
%\frametitle{Données et indices numériques, Virulence}

%\begin{deft}[Virulence]
%\begin{small}
%Soit $v$ un virus, on définit
%\begin{enumerate}
%\item L'indice infectieux $I^0_v$ d'un virus $v$ mesure le risque \enlight{a priori}.
%$$I^0_v = \frac{\mbox{Nombre de fichiers infectables par } v }{\mbox{Nombre total de fichiers du système}}$$
%\item L'indice infectieux $I^1_v$ d'un virus $v$ mesure le risque \enlight{a posteriori}.
%$$I^1_v = \frac{\mbox{Nombre de fichiers infectés par } v }{\mbox{Nombre  de fichiers infectables par } v}$$
%\item La virulence $V_v$ d'un virus $v$ est définit comme
%$$V_v = I^0_v \times I^1_v = \frac{\mbox{Nombre total de fichiers du système}}{\mbox{Nombre de fichiers infectés par } v }$$
%\end{enumerate}
%\end{small}
%\end{deft}
%\end{frame}


%\begin{frame}
%\frametitle{Données et indices numériques, Virulence}
%\begin{itemize}
%\item Fichiers considérés : uniquement fichiers exécutables ou fichiers documents (ps, pdf, word,\ldots).
%\item Ces indices ont un sens que dans un environnement bien défini (OS,etc.)
%\item Remarquons que les indices $I_v^0,I_v^1$ et $V_v$ ne considère que le risque d'infection et ne considère pas du t%out la charge finale.
%\item Ces indices sont relativement faciles à obtenir.
%\item Dans le cas des vers, les indices précédents sont définis en considérant non pas des fichiers mais des machines i%nfectées.
%\end{itemize}
%\end{frame}






%\begin{frame}
%\frametitle{Données et indices numériques - Détectabilité}

%\begin{small}
%\begin{proposition}[degré de détectabilité]
%Le degré de détectabilité d'une infection informatique  est inversement proportionnel à la durée de la phase d'incubation et proportionnel au nombre d'infections survenues dans ce système. Autrement dit
%$$\mbox{Détectabilité} = C \times \frac{\mbox{ Nombre de copies }}{\mbox{ Temps incubation }}$$
%où $0 \leq C \leq 1$. 
%\end{proposition}
%\begin{itemize}
%\item La constante $C$ décrit un certain nombre de paramètre : qualité de programmation, technique anti-antivirale, nat%ure de la charge finale.
%\item Cela permet donc de prendre en compte la totalité de l'activité du virus.
%\item Le degré de risque d'une infection informatique varie grosso modo inversement avec le degré de détectabilité de c%ette infection.
%\end{itemize}
%\end{small}
%\end{frame}



\section{Mode d'action des virus }

\subsection{Virus par écrasement de code}


\begin{frame}
\frametitle{Virus par écrasement de code}
Lorsque le virus est exécuté : 
\begin{itemize}
\item La routine de recherche identifie des cibles.
\item L'infection se fait simplement en écrasant le code cible par son propre code.
\item Ces virus sont en général dépourvu de charge finale, car le  mode d'infection du virus en est une à elle seule.
\end{itemize}
\end{frame}

\begin{frame}
\frametitle{Virus par écrasement de code}

Trois cas d'écrasement sont possibles
\medskip

\enlight{\ding{202}} Le virus écrase la partie initiale du code de la cible :

%\begin{itemize}
%\item L'en-tête spécifique de l'exécutable dont la fonction est de structurer les données est absent.
%\item Le code afin de faciliter la projection en mémoire par l'OS est alors absent.
%\item Le programme infecté ne pourra pas se lancer.
%\end{itemize}

\medskip


\enlight{\ding{203}} Le virus écrase le code cible en partie centrale ou terminale. 
%\begin{itemize}
%\item Le virus doit alors installer une fonction de saut vers l'adresse de son %propre code en début du programme infecté s'il veut la main avant le programme cible.
%\item Selon les cas le programme cible ne se lancera pas et/ou son exécution avortera en cours de processus.
%\item  Dans ce cas on obtient une dose de furtivité : le dysfonctionnement est imputé par l'utilisateur à un boggue du programme plutôt qu'à un virus.
%\end{itemize}

%\end{frame}



%\begin{frame}
%\frametitle{Virus par écrasement de code}
\medskip

\enlight{\ding{204}}  Le virus remplacement simplement le code cible par son propre code. %Technique peu courante et facilement détectable puisque tous les exécutables infectés ont la même taille, et même code.

\begin{center}
\includegraphics[width=8cm]{ecrasement}
\end{center}

\end{frame}





\subsection{Virus par recouvrement de code (1/2)}

\begin{frame}

\frametitle{Virus par recouvrement de code}

Les virus fonctionnant selon ce mode vont accoler leur code à celui de la cible.
Deux cas:
%Il en résulte une augmentation de la taille du programme infecté.


%Le virus a deux possibilités pour accoler son code :
\begin{itemize}
\item  \enlight{En tête du code de la cible.}% (type \textit{prepend} ou ajout en position initiale). C'est le cas le moins fréquemment réalisé car le plus compliqué. 
C'est compliqué car on doit modifier dans l'entete :% entre autre données
\begin{itemize}
\item Longueur du fichier et éventuellement des sections de code.
\item Re-calcul des adresses.
\item En plus un \enlight{saut} vers le début du code cible doit être ajouté. % pour que le code cible s'exécute.
\end{itemize}

\end{itemize}
\end{frame}

\begin{frame}
\frametitle{Virus par recouvrement de code (2/2)}
\begin{itemize}
\item \enlight{en fin du code de la cible.} 
\begin{itemize}
\item Le début du code cible doit alors être remplacé par un saut vers le code viral.
\item Le code remplacé doit être exécuté à la fin du code viral.
\item Suivit d'un saut vers le code cible doit être effectué.
\end{itemize}

\begin{center}
\includegraphics[width=10cm]{recouvrement}
\end{center}


\end{itemize}


\end{frame}

\subsection{Virus par entrelacement de code}

\begin{frame}
\frametitle{Virus par  entrelacement  de code}
Ces virus exploitent essentiellement le format PE des exécutables 32 bits de Windows. Cette entête permet lors de la projection du code en mémoire :
\begin{itemize}
\item de donner les informations adéquates pour l'installation en mémoire (établissement d'une image mémoire);
\item de permettre la mise en commun optimale pour plusieurs processus de fichier EXE et DLL.
\end{itemize}
Ce format se prête particulièrement bien  à l'écriture de virus.
\end{frame}



%\begin{frame}
%\frametitle{Virus par  entrelacement  de code - fichier PE}
%Un fichier PE se compose
%\begin{itemize}
%\item D'une entête DOS qui permet de lancer le programme sous DOS et d'afficher le message indiquant que l'application ne fonctionne que sous Windows.
%\item De l'entête PE proprement dite. Il comprend deux structures de données essentielles, renseignées lors de la compilation et de l'édition dynamique et indispensables au lancement correct de l'application.
%\end{itemize}

%\end{frame}



%\begin{frame}
%\begin{columns}[c]
%\begin{column}{4cm} %{0.5\textwidth}
%\only<1>{\includegraphics[width=4cm]{fichierPE}}
%\only<2>{\includegraphics[width=4cm]{fichierPE-fileheader}}
%\only<3>{\includegraphics[width=4cm]{fichierPE-optionalheader}}
%\only<4>{\includegraphics[width=4cm]{fichierPE-sectionheader}}
%\includegraphics[width=4cm]{fichierPE-sectionheader}
%\end{column}

%\begin{column}[c]{6cm} 

%\only<2>{
% l'\texttt{IMAGE\_FILE\_HEADER} est 
%\smallskip


%\color{red}
%\begin{tt}
%\noindent typedef struct\{\\
%WORD Machine;\\
%WORD NumberOfSection;\\
%DWORD TimeDateStamp;\\
%DWORD PointerToSymbolTable;\\
%DWORD NumberOfSymbols;\\
%Word SizeOfOptionHeader;\\
%Word Characteristics;\\
%\} IMAGE\_FILE\_HEADER\\
%\end{tt}
%}

%\only<3>{\color{red}

%\begin{tt}

%\begin{small}
%\noindent typedef struct\{\\
%.... \\
%DWORD SizeOfCode;\\
%DWORD SizeOfInitializedData;\\
%DWORD SizeOfUnInitializedData;\\
%DWORD AddresOfEntryPoint;\\
%DWORD BaseOfCode;\\
%DOWRD BaseOfData;\\
%DWORD ImageBase \slash\slash adresse \\
%\qquad\qquad  \slash\slash de chargement\\
%DWORD SectionAlignment;\\
%....\\
%\} IMAGE\_OPTIONAL\_HEADER;
%\end{small}

%\end{tt}
%}


%\only<4>{
%Seule les \texttt{NumberOfRvaAndSizes} premières entrées sont renseignées
%\begin{tt}{
%\color{red}
%typedef struct\{\\
%DWORD VirtualAddress; \slash\slash Adresse RVA du
%\hspace{2cm}  \slash\slash  debut de la section
%DWORD Size; \slash\slash taille de la section \\
%\} IMAGE\_DATA\_DIRECTORY;
%}
%\bigskip


%Les autres valent 0.
%\end{tt}
%}

%\end{verbatim}
%\end{column}


%\end{columns}
%\end{frame}

%\begin{frame}
%\frametitle{Virus par  entrelacement  de code}

%\begin{itemize}
%\item Toutes les adresses contenues dans le format PE sont des adresses relatives (RVA)
%\item Lors de la projection en mémoire, l'adresse mémoire est obtenue en calculant
%$$\texttt{RVA + ImageBase}.$$
%\item Les taille mémoire des sections ont une taille mémoire multiplie de 512 octets.
%\item L'infection utilise les octets libre de chaque section.
%\item Il n'y a pas d'augmentation de taille mémoire lors de l'infection.
%\end{itemize}

%\end{frame}


\begin{frame}
\begin{center}
\includegraphics[width=10cm]{entrelacement}
\end{center}

\end{frame}

\subsection{Virus par accompagnement de code}

\begin{frame}
\frametitle{Virus par accompagnement de code}
\begin{itemize}
\item Méthode la moins connue.
\item Le code cible n'est pas altéré.
\item le code viral prend le nom du code source, et exécute le code cible avec la commande \texttt{exec} ou \texttt{system}
\end{itemize}
\begin{center}
\includegraphics[width=10cm]{accompagnement}
\end{center}

\end{frame}

\begin{frame}
\frametitle{Accompagnement de code - Exemple sous Linux}

\begin{itemize}
\item Sous linux, le virus modifie le PATH en ajoutant en tête le répertoire courant ``.''
\item Dans le répertoire courant, une copie du virus est faite sous le nom de \texttt{ls}
\item Le code du virus est modifié afin qu'il lance à la fin de son exécution la commande \texttt{ls} du système.
\end{itemize}

Une approche similaire peut être faite sous Windows.
\end{frame}

\subsection{Virus de code source}

\begin{frame}
\frametitle{Virus de code source}
\begin{itemize}
\item La cible ici est le code source d'un programme C.
\item Le programme ainsi infecté doit être recompilé afin de produire un exécutable valide.
\item Une rapide analyse du code infecté pourrait trahir la présence du virus, il faut donc dupliquer le code de façon plus évolué.
\item L'intérêt d'une infection par code source vient du fait que le binaire obtenu est totalement homogène.
\end{itemize}
\end{frame}



\begin{frame}
\frametitle{Virus de code source}
\begin{center}
\includegraphics[width=8cm]{viruscodesource}
\end{center}
\end{frame}


\subsection{Les techniques anti-antivirales}


%\begin{frame}
%\frametitle{Les techniques anti-antivirales}
%\begin{deft}[Sécurité]
%Ensemble  des mesures et techniques destinées à contrer les actions malveillantes contre un système, actions dont la nature propre est s'adapter aux protections qui lui sont opposées.
%\end{deft}
%\end{frame}

\begin{frame}
\frametitle{Les techniques anti-antivirales}
\begin{itemize}
\item %Dans le cadre de la lutte antivirale, il est alors parfaitement logique que les virus, vers ou autres infections informatiques,
Les malware  évoluent de facon contrer et à défaire les protections mises par les anti-virus ou les pare feux.
\medskip

\item  \enlight{La furtivité.}  Sert à leurrer l'utilisateur, le système et les logiciels de protection en faisant croire à l'absence d'un code malveillant. % Par exemple en dissimulant le code comme dans l'infection par entrelacement.
\medskip

\item \enlight{Le polymorphisme.} Les anti-virus fonctionnant en grande partie sur la recherche de \alert{ signatures virales}, le but du polymorphisme est de faire varier, de copie en copie virale le code viral. %tout élément fixe pouvant être exploité par l'antivirus pour identifier le virus. 
\end{itemize}
\end{frame}

\begin{frame}
\frametitle{Les techniques anti-antivirales : polymorphisme}

Les deux principales techniques sont les suivantes


\begin{itemize}
\item Réécriture de code par utilisation de code équivalent. Par exemple en langage C
\begin{flushleft}
\begin{tt}
if( flag ) infection();
else charge\_finale();
\end{tt}
\end{flushleft}
pourra être réécrite comme suit
\begin{flushleft}
\begin{tt}
(flag) ? infection() : charge\_finale();
\end{tt}
\end{flushleft}
\medskip

\item Utilisation de technique de chiffrement basique sur tout ou partie du virus ou du ver. %En général, il s'agit plus d'un procédé de masquage (XOR bit à bit) avec une valeur constante.

%Le code viral commence par une procédure en clair dont le but est de déchiffrer le reste du code. Il est alors nécessaire de changer à chaque infection de procédure de chiffrement.

\end{itemize}
\end{frame}



\section{Classification des virus et des vers}


\subsection{Classification des virus}
\begin{frame}
\frametitle{Classification des virus}
On peut classer les virus et les vers suivants différents critères
\begin{itemize}
\item selon le format visé : virus {\color{blue} d'exécutables } ou de {\color{blue} documents.}
\item selon l'organe cible : virus de {\color{blue} secteur de boot, de pilotes de périphériques},\ldots.
\item selon le langage de programmation : virus en assembleur, virus de code source, virus en script.
\item selon le comportement : virus blindés, virus lents ou rapides, rétrovirus, virus résidents, virus polymorphes ou furtifs,\ldots
\end{itemize}
\end{frame}



%\begin{frame}
%\frametitle{Virus d'exécutable}
%\begin{itemize}
%\item L'infection concerne une cible binaire à partir d'un fichier binaire infecté. 
%\item Il s'agit en général de mécanisme de bas niveau qui obligent, en règle générale, à utiliser l'assembleur.
%\end{itemize}

%Les cibles principales sont les suivantes
%\begin{itemize}
%\item Les fichier \texttt{*.COM} qui sont de taille $<  64KO$, donc un seul segment mémoire. La contrainte forte lors de l'infection est de garder un binaire $< 64KO$.
%\item Exécutable \texttt{*.EXE}, on a vu les stratégies possibles dans la section précédentes (entrelacement, accompagnement, etc.).
%\item Fichier exécutables au format ELF (Unix). 
%\end{itemize}
%
%C'est le format de ces fichiers qui va dicter le mécanisme d'action du virus.
%\end{frame}


\begin{frame}
\frametitle{Virus de document ou macro virus.}
Nous adopterons la définition suivante :

\begin{deft}[Virus de documents]
Un virus de document est
\begin{itemize}
\item Un code viral contenu dans un {\color{blue} fichier de données, non exécutable.}
\pause
\item Il est activé par {\color{blue} un  interpréteur}  contenu de fa\c con native dans l'application associée au format de ce fichier (déterminé par son extension). 
\pause
\item L'activation du code malveillant est réalisée
\begin{itemize}
\item soit par une fonctionnalité prévue dans l'application, 
\item soit en vertu d'une faille interne de l'application considérée.
\end{itemize}
\end{itemize}
\end{deft}

\end{frame}


%\begin{frame}
%\frametitle{Virus de document ou macro virus.}
%On évalue le risque lié à l'infection d'un document lié au condition d'exécution du code du document. 
%\begin{enumerate}
%\item Le format de fichier contient toujours du code, qui est directement exécuté à l'ouverture du fichier.
%\item Le format contient parfois du code, qui  peut s'exécuter directement.
%\item Le format contient parfois du code qui ne s'exécuter qu'après confirmation de l'utilisateur.
%·\item  Le format contient parfois du code qui ne peut s'exécuter qu'après une action volontaire de l'utilisateur.
%\item Le format ne peut jamais contenir le code.
%\end{enumerate}
%\end{frame}



\begin{frame}
\frametitle{Virus de documents : risques associés aux formats connus}
\begin{center}
\begin{small}
\begin{tabular}{|c|c|c|c|}
\hline
Format & Extensions de fichier & Risque & Type \\
\hline
\hline
Script WSH & VBS,JS,VBE,JSE,WSF,WSH & 1 & texte \\
\hline
Word & DOC,DOT,WBK,DOCHTML & 2,3, & binaire \\
\hline 
Excel & XLS,XL?,SLK,XLSHTML & 2,3 & binaire \\
\hline
Powerpoint & PPT, POT,PPS, & 2,3 & binaire \\
 & PPA,PWZ,PPTHTML, POTHTML & & \\
\hline
Acces & MDB,MD?,MA?,MDBHTML & 1 &binaire \\
\hline
RTF & TRTF & 4 & texte \\
\hline
Shel Scrap & SHS & 1 & binaire \\
\hline
HTML & HTML,HTM,\ldots & 2 & texte\\
\hline
XHTML & XHTML,XHT,\ldots & 2 & texte\\
\hline
Adobe Acrobat & PDF & 2 & texte \\
\hline
PostScript & PS & 1,2 & texte \\
\hline
tex,latex & TEX & ? & texte \\
\hline
\end{tabular}
\end{small}
\end{center}
\end{frame}



%\begin{frame}
%\frametitle{Virus de document - Statistique}

%\begin{itemize}
%\item Les formats les plus touché par les infections sont ceux de la suite Office de Microsoft.
%\item Cela s'explique en partie par la popularité de ces logiciels et formats de fichiers.
%\item Statistique (ci-dessous) des attaques sur les différentes version de la suite office.
%\end{itemize}


%\begin{center}
%\includegraphics[width=7cm]{stat-macro-virus}
%\end{center}

%\end{frame}



\begin{frame}
\frametitle{Virus de Bios - Virus de démarrage}
\noindent  {\color{blue} Virus de Bios :}
\begin{itemize}
\item Lors de la mise sous tension de l'ordinateur, le code contenu dans le Bios est activé. 
\item Son rôle est de tester l'environnement matériel et ensuite de passer le contrôle au secteur de démarrage.
\item Un virus de Bios modifie le Code BIOS.
\end{itemize}
\noindent {\color{blue} Virus de démarrage :}
\begin{itemize}
\item le code viral est placé dans le secteur de démarrage du disque dur.
  \item Il s'exécute avant tout autre programme, donc ca lui permet de désactiver la lutte antivirale avant qu'elle s'exécute.
  \end{itemize}
\end{frame}

%\begin{frame}
%\frametitle{Virus de Bios}
%Un virus BIOS présente plusieurs avantages fondamentaux
%\begin{itemize}
%\item il n'est détecté par aucun antivirus. Ces derniers surveillent au niveau du BIOS le secteur de démarrage maître et au niveau du système d'exploitation le disque dur et l'activité système. Aucun contrôle du BIOS lui même n'est effectué.

%\item du fait de son antériorité de lancement, tout programme situé au niveau du Bios ou lancé par lui peut agir sur les données présentes sur la disque.
%\item  Tout programme lancé par le Bios accède à n'importe quelle donnée sur le disque sans limitation. Les possibilités en terme d'infection et de charges finales sont sans limite.

%\end{itemize}
%\end{frame}

%\begin{frame}
%\frametitle{Virus de démarrage}
%Le virus se copie dans le secteur de démarrage : 
%\begin{itemize}
%\item Le secteur de démarrage maître (Master Boot Record ou MBR) prend le relais du Bios. 
%\item Son rôle est de lancer l'un des systèmes d'exploitation présents sur le disque. 
%\item Ce secteur est un exécutable de 512 octets.
%\end{itemize}
%L'intérêt du virus de démarrage est de s'exécuter avant tout autre programme
%Comme dans le cas du virus de BIOS, l'intérêt du virus de démarrage est de s'exécuter avant tout autre programme (hors BIOS) dont un éventuel antivirus. 

%\end{frame}

%\begin{frame}
%\frametitle{Virus de démarrage}
%Deux techniques d'infection du secteur de démarrage existent
%\begin{enumerate}%
%\item Le virus est en réalité un secteur de boot viral, 
%\begin{itemize}
%\item Il écrase et remplace le secteur original sain. 
%\item Il possède toutes les fonctionnalités d'un programme de démarrage en y incluant les capacités virales. 

%\item Cela limite les fonctionnalité du virus (faible taille du code).
%\end{itemize}

%\item Virus plus complexes, possédant des fonctionnalités plus évoluées (furtivité, charge finale) la taille du virus dépasse largement la taille d'un secteur de démarrage

%\begin{itemize}
%\item Le secteur de démarrage viral SDV qui va venir remplacer le secteur de démarrage original sain.
%\item Plusieurs secteurs additionnels viraux. Ces secteurs sont généralement dissimulés ou chiffré. C'est le secteur SDV qui rassemblera ces secteur disséminé lors de l'exécution.
%\item Une copie du secteur de démarrage sain. Le virus une fois l'installation en mémoire effectuée pour pouvoir infecter d'autres secteurs de démarrage redonne le contrôle au secteur original sain qui  lancera le système d'exploitation. 
%\end{itemize}

%\end{enumerate}
%\end{frame}


\begin{frame}
\frametitle{Virus résident}
C'est un virus {\color{blue} constament en cours d'exécution. }
\begin{itemize}
\item Un virus résident est un virus qui une fois exécuté reste actif: seul l'arrêt de la machine ou un \texttt{kill} peut l'arrêter.
\item Il utilise des technique diverse pour pouvoir être ré-exécuter après un arrêt de la machine (lancement au démarrage, etc..).
\item Il peut interférer avec le système (court-circuiter certaines intéruptions systèmes par exemple)
\end{itemize}
\end{frame}


%\end{document}

%\begin{frame}
%\frametitle{Virus lents - Virus rapides}


%Ces virus en fait sont le plus souvent de simples virus d'exécutables, dont l'infection utilise une des techniques de la section précédente.

%\begin{itemize}
%\item les virus lents, en mode résident, n'infectent que les fichiers exécutables qui sont modifiés ou créés.  Le but est de leurrer les antivirus, notamment ceux fonctionnant par contrôle d'intégrité en faisant passer la modification du code et l'action infectieuse proprement dite comme légitimes.
%\item Les virus rapides au contraire eux aussi résidents infectent cette fois les fichiers qui sont exécutés ou ouverts. Cet événement est fréquent ( d'où le nom de virus rapides), tout particulièrement lors de l'exécution d'un antivirus (ouverture de fichier, recherche de signature et émulation de code). Le virus cette fois emboîte le pas à l'antivirus, et ce dernier ne voit que sa propre action. 
%\end{itemize}
%\end{frame}




\begin{frame}
\frametitle{Virus psychologique (hoax - cannular)}

%C'est un phénomène relativement récents et très particulier: c'est une sorte de duplication de code social.
C'est une sorte de code viral qui s'exécute dans le cerveau.
\bigskip

\begin{definition}[Virus psychologique]\begin{itshape}
Un \enlight{virus psychologique} est une désinformation incitant l'utilisateur, par des techniques d'ingénieurie sociale, a produire des effets équivalents à celui d'un virus ou d'un ver : propagation et action offensive. 
\end{itshape}
\end{definition}

\bigskip
{\color{blue} Exemple:} circulation d'un mail avertissant d'une infection, solutionner par une suppression de fichier système.

\end{frame}




\subsection{Classification des vers}


\begin{frame}
  \frametitle{Classification}
 \noindent {\color{blue} Ver Simple :}
\begin{itemize}
\item Ce sont eux qui légitimement peuvent s'appeler ver.
\item Ils exploitent en général des failles logicielles et réseau permettant l'exécution de programmes sur une machine distante.
%\item Ils exploitent aussi des faiblesses dans les protocoles réseaux (mots de passe faible, authentification sur l'IP,\ldots).
\end{itemize}
\noindent {\color{blue} Macro vers  (hybrides virus et vers) :}
\begin{itemize}
\item Le mode de dissémination se fait pas des pièces jointes contenant des documents bureautiques infectés aux utilisateurs du carnet d'adresse de l'utilisateur.
\end{itemize}
\noindent {\color{blue} Vers d'emails :}
\begin{itemize}
\item Propagation par pièce jointe contenant un code malicieux.
\item Le code est activité soit par l'utilisateur soit par l'application de courrier électronique.
  \item L'exemple le plus célèbre est le ver I LOVE YOU.
  \end{itemize}
\end{frame}

%\end{document}

%\section{La lutte antivirale}


\section{Conclusion}

\begin{frame}
  \frametitle{Que retenir ?}
  C'est un chapitre donnant des informations générales sur les logiciels malveillants : 
  \begin{itemize}
  \item Connaître les différents type de logiciels malveillants.
  \item Distinguer les virus et les vers.
  \item Connaître les différents modes de copie des virus.
    \item Connaître les différentes formes de propagation des virus.
    \end{itemize}
\end{frame}
\end{document}
